% ============================================================
%  APPENDIX TO CHAPTER 4 -- EXERCISES S4.1-S4.25
%  Imported from Tao I/II, ECNU, USTC, Yu Pin.  Not Rudin's.
% ============================================================
\ssection{Appendix exercises}

\begin{roadmap}[How these differ from Rudin's twenty-six]
Rudin's own twenty-six exercises, which build \emph{his} machinery --- the
uniform-continuity chain, the $\rho_E$ toolbox, the counterexample gallery
--- end the chapter proper. These twenty-five come from the four comparison
texts and fill the other half of the subject.

\medskip
\raggedright
\textbf{Group I (1--6): limits and asymptotics.} Computational problems in the
$o$/$\sim$ language of Supplement B, plus three rigidity problems --- a
periodic function that dies at infinity is dead everywhere, and its
relatives. Nothing in Rudin's set resembles these.

\smallskip
\textbf{Group II (7--11): monotonicity, inverses, functional equations.}
Where order replaces compactness. Exercise 9 (Cauchy's functional equation)
and 10 are the standard route from ``continuous at one point'' to
``continuous, and linear.''

\smallskip
\textbf{Group III (12--16): what the intermediate value theorem really
gives.} The discrete mean value property, the universal chord theorem, and a
contraction mapping five chapters before Rudin needs one (9.23).

\smallskip
\textbf{Group IV (17--21): uniform continuity as a working tool.} Prove the
criteria of Supplement E, then use them. Compare with Rudin's 4.8--4.13,
which prove the same facts one at a time by hand.

\smallskip
\textbf{Group V (22--25): structure.} Oscillation, prescribed discontinuity
sets, path-connectedness, and one spectacular payoff of ``continuous plus
dense determines everything'' --- a proof of the Cayley--Hamilton theorem
that uses no algebra beyond the diagonalisable case.

\medskip
\textbf{Marked $\star$:} harder, or requiring a result from a later chapter.
\end{roadmap}

\begin{enumerate}[label=\textbf{S4.\arabic*.}, leftmargin=3.4em, itemsep=9pt,
                  topsep=8pt, labelsep=0.5em]

% ---------------- Group I ----------------
\item Let $f$ be periodic on $\R^1$ with $\lim_{x\to+\infty} f(x) = 0$.
Prove that $f$ is identically zero. No continuity is
assumed.\kn{Fix $x$ and push it to infinity by its own period: $f(x) =
f(x+nT)$ for every $n$, and the right side tends to $0$. The whole content
is that periodicity lets one evaluate a function arbitrarily far out
\emph{without moving it}. Compare Rudin's Exercise 4.4: there density does
the transporting, here periodicity does.}
\provenance{ECNU \cjk{习题}{Ex} 3.3 \#7; USTC Ch.\ 1 suppl.\ \#13.}

\item Suppose $f$ is defined on $(0,+\infty)$, satisfies $f(x^2) = f(x)$ for
all $x$, and $\lim_{x\to0^+} f(x) = \lim_{x\to+\infty} f(x) = f(1)$. Prove
that $f$ is constant.\kn{Iterate the hypothesis backwards: $f(x) =
f(x^{1/2}) = f(x^{1/4}) = \cdots = f(x^{1/2^n})$, and $x^{1/2^n} \to 1$ for
every $x>0$. The two limit hypotheses are there only to handle $x$ where
the iteration escapes. This is the shape of every functional-equation
rigidity argument: find the orbit, find its limit point, evaluate there.}
\provenance{ECNU \cjk{总练习题}{review} 3 \#13.}

\item $\star$~Let $f$ be defined on $(a,+\infty)$ and bounded on every finite interval
$(a,b)$, and suppose $\lim_{x\to+\infty}\bigl[f(x+1)-f(x)\bigr] = A$. Prove
that $\lim_{x\to+\infty} f(x)/x = A$.\kn{The continuous Stolz--Ces\`aro
theorem. Write $f(x) = f(x_0) + \sum_{k<n}\bigl[f(x_0+k+1)-f(x_0+k)\bigr]$
where $x = x_0 + n$ with $x_0$ in a fixed unit window, then average: the
increments are eventually near $A$, and the finitely many early ones are
diluted by $1/x$. It is 3.20's averaging idea run on a continuous variable.}
\provenance{ECNU \cjk{总练习题}{review} 3 \#14; Yu Pin, HW5 A9$^*$ (p.\,127),
where it is attributed to Cauchy rather than to Stolz.}

\item Show that $f(x) = x\sin x$ is unbounded on $(0,+\infty)$ but is
\emph{not} an infinitely large quantity as $x \to +\infty$; that is,
$\lim_{x\to+\infty} f(x) \ne \infty$.\kn{Two sequences settle it: $x_n =
2n\pi$ gives $f = 0$, while $x_n = 2n\pi + \pi/2$ gives $f \to +\infty$.
The point of Supplement B's S4.8 --- ``unbounded'' is a statement about the
range, ``tends to infinity'' is a statement about the tail, and only the
second forbids returning.}
\provenance{ECNU bk.\ p.\,60; USTC Ex.\ 1.2.12.}

\item Using equivalent infinitesimals, evaluate
$\displaystyle\lim_{x\to0}\frac{\tan x - \sin x}{\sin^3 x}$ and
$\displaystyle\lim_{x\to0}\frac{\sqrt{1+x^2}-1}{1-\cos x}$. Then explain
precisely why substituting $\tan x \sim x$ and $\sin x \sim x$ into the
numerator of the first gives the wrong answer $0$.\kn{Answers $\tfrac12$ and
$1$. The explanation is the content of the pitfall at S4.7: each equivalence
is accurate to first order, the subtraction annihilates exactly the first
order, and what survives lives at third order where neither equivalence says
anything. Factor first --- $\tan x - \sin x = \sin x(1-\cos x)/\cos x$ ---
and every factor may then be replaced.}
\provenance{ECNU \cjk{例}{Ex} 2 and note, bk.\ p.\,59.}

\item Find all asymptotes of the curve $y = x^3/(x^2+2x-3)$, and prove that
your oblique asymptote satisfies the definition (the distance from
$\bigl(x,f(x)\bigr)$ to the line tends to $0$).\kn{Long division is the
whole computation; the definition check is the observation that the distance
to $y=ax+b$ is $|f(x)-ax-b|/\sqrt{1+a^2}$, a fixed multiple of the vertical
gap. Worth doing once to see that ``vertical gap'' and ``perpendicular
distance'' differ only by that constant --- which is why the limits defining
$a$ and $b$ may be taken vertically.}
\provenance{ECNU \cjk{例}{Ex} 6, bk.\ pp.\,61--62.}

% ---------------- Group II ----------------
\item Let $f$ be continuous and one-to-one on an interval $I$. Prove that
$f$ is strictly monotone on $I$. Deduce that a continuous injection of
$[a,b]$ into $\R^1$ has a continuous inverse, without invoking
compactness.\kn{Suppose not: then some $u<v<w$ has $f(v)$ outside the
interval with endpoints $f(u)$ and $f(w)$. Apply the intermediate value
theorem on $[u,v]$ and on $[v,w]$ to produce two points with the same value.
The deduction is S4.13. Contrast Rudin 4.17, which gets the same conclusion
from compactness and works in any metric space --- neither theorem contains
the other.}
\provenance{Tao I, Ex.\ 9.8.3; USTC \cjk{定理}{Thm} 3; Yu Pin, Thm 54.}

\item Prove that there exists a function $f : \R^1 \to \R^1$ that is
discontinuous at every point while $|f|$ is continuous at every
point.\kn{Modify Dirichlet: $f = 1$ on the rationals, $-1$ elsewhere. Then
$|f| \equiv 1$. The lesson is that $|\cdot|$ destroys information --- the
implication ``$f$ continuous $\Rightarrow |f|$ continuous'' has no converse,
and every proof that passes to absolute values gives something away.}
\provenance{USTC Ex.\ 2.1.5.}

\item Suppose $f : \R^1 \to \R^1$ satisfies $f(x+y) = f(x) + f(y)$ for all
$x,y$, and is continuous at the single point $0$. Prove that $f$ is
continuous on all of $\R^1$ and that $f(x) = f(1)\,x$.\kn{Cauchy's
functional equation. Continuity everywhere follows because $f(x+h)-f(x) =
f(h)$ --- the equation makes the increment independent of position, so
continuity at one point is continuity everywhere. Then the equation forces
$f(q) = f(1)q$ for rational $q$ by induction, and Rudin's Exercise 4.4
(agreement on a dense set) finishes it. Without \emph{some} regularity the
conclusion is false, but the counterexamples need a Hamel basis and hence
the axiom of choice.}
\provenance{ECNU \cjk{总练习题}{review} 4 \#9.}

\item Let $f$ be defined on $\R^1$, continuous at the two points $0$ and
$1$ only (as far as is assumed), and satisfy $f(x^2) = f(x)$ for all $x$.
Prove that $f$ is constant.\kn{Sharper than S4.2: continuity at exactly two
points suffices. Note $f$ is even, so it is enough to work on $x \ge 0$;
for $x>0$ the iteration $x^{1/2^n} \to 1$ uses continuity at $1$, and
$f(0) = \lim_{x\to0}f(x)$ uses continuity at $0$. Observe which points the
iteration funnels into --- those are exactly the points where continuity is
required.}
\provenance{ECNU \cjk{总练习题}{review} 4 \#10.}

\item Determine, with proof, which of the following are possible for a
continuous real function: domain $[0,1]$ with range $(0,\infty)$; domain
$[0,1]$ with range $(0,1)$; domain $[0,1]$ with range
$[0,1]\cup[2,4]$; domain $(0,1)$ with range $[2,\infty)$.\kn{Four probes of
the same two theorems. The first three are impossible --- 4.15 makes the
image compact, hence closed and bounded, and 4.22 makes it connected, hence
an interval. The fourth is possible; construct one. A good check on whether
you know what compactness and connectedness each contribute.}
\provenance{USTC Ex.\ 2.2.10.}

% ---------------- Group III ----------------
\item Let $f$ be continuous on $[a,b]$ and let $x_1,\dots,x_n \in [a,b]$.
Prove that there is a $\xi \in [a,b]$ with
\[ f(\xi) = \frac{f(x_1)+\cdots+f(x_n)}{n} . \]
More generally, prove the same for any weighted average
$\sum q_i f(x_i)$ with $q_i > 0$ and $\sum q_i = 1$.\kn{The average lies
between $\min_i f(x_i)$ and $\max_i f(x_i)$, both of which are attained
values of $f$; now apply the intermediate value theorem between the two
points where they are attained. This is the discrete ancestor of the mean
value theorem for integrals (6.x) and of the notion of a barycentre.}
\provenance{ECNU \cjk{习题}{Ex} 4.2 \#19; USTC Ex.\ 2.2.7.}

\item \emph{(Universal chord theorem.)} Let $f$ be continuous on $[0,1]$
with $f(0) = f(1)$. Prove that for every positive integer $n$ there is a
$\xi \in [0,1-1/n]$ with
\[ f\Bigl(\xi + \frac1n\Bigr) = f(\xi). \]
Show by example that the conclusion can fail for a chord length that is not
of the form $1/n$.\kn{Put $F(x) = f(x+1/n) - f(x)$ on $[0,1-1/n]$ and sum:
$F(0) + F(1/n) + \cdots + F\bigl((n-1)/n\bigr) = f(1)-f(0) = 0$. A finite
list of reals summing to zero either is all zeros or contains both signs, so
$F$ has a zero by 4.23. The telescoping sum is the whole idea, and it is
why only reciprocal lengths work.}
\provenance{ECNU \cjk{总练习题}{review} 4 \#8; USTC Ch.\ 2 suppl.\ \#5.}

\item Prove that for every positive integer $n$ the equation
\[ x^n + x^{n-1} + \cdots + x = 1 \]
has exactly one positive root $x_n$; prove that $(x_n)$ converges, and find
its limit.\kn{Existence and uniqueness from 4.23 plus strict increase of the
left side on $(0,\infty)$. For the limit, note $x_n \in (\tfrac12,1)$ and
that the sum telescopes: $x(1-x^n)/(1-x) = 1$. Show $(x_n)$ is decreasing and
bounded below, then pass to the limit in the closed form to get $\tfrac12$.
A clean synthesis of Chapters 3 and 4.}
\provenance{USTC Ch.\ 2 suppl.\ \#9.}

\item \emph{(Contraction mapping on an interval.)} Let $f$ map $[a,b]$ into
$[a,b]$ and suppose there is a constant $k$ with $0<k<1$ such that
\[ |f(x)-f(y)| \le k|x-y| \qquad (x,y \in [a,b]). \]
Prove that $f$ has a unique fixed point $x_0$, and that for every starting
value $x_1 \in [a,b]$ the sequence $x_{n+1} = f(x_n)$ converges to
$x_0$.\kn{Uniqueness is immediate from the inequality. For existence, do not
use 4.23 --- use the iteration: $|x_{n+1}-x_n| \le k^{n-1}|x_2-x_1|$ makes
$(x_n)$ Cauchy by comparison with a geometric series (3.26), and 3.11 gives
the limit, which is a fixed point by continuity. This is Rudin's 9.23 in one
dimension, and the proof is identical --- worth doing now so that 9.23 is
recognition rather than novelty. Compare Rudin's Exercise 4.14, which gets a
fixed point from the intermediate value theorem alone but no iteration.}
\provenance{USTC Ch.\ 2 suppl.\ \#8.}

\item $\star$~\emph{(Croft.)} Let $f$ be continuous on $\R^1$ and suppose that for every
$\delta > 0$,
\[ \lim_{n\to\infty} f(n\delta) = 0 . \]
Prove that $\lim_{x\to+\infty} f(x) = 0$.\kn{Surprisingly hard, and the
tool is Baire (Rudin's Exercise 3.22). Suppose the conclusion fails; build
closed sets $F_N = \{\delta \ge 1 : |f(n\delta)| \le \varepsilon \text{ for
all } n \ge N\}$, whose union is $[1,\infty)$, so one has interior; then the
dilates $\{n\delta : \delta \in I,\ n \ge N\}$ of an interval $I$ eventually
cover a half-line. The moral: a hypothesis quantified over \emph{every}
$\delta$ is a category hypothesis in disguise.}
\provenance{Yu Pin, HW5 problem G (p.\,129).}

% ---------------- Group IV ----------------
\item Prove the sequential criterion S4.15: $f$ is uniformly continuous on
$E$ if and only if $d\bigl(f(p_n),f(q_n)\bigr) \to 0$ for every pair of
sequences in $E$ with $d(p_n,q_n) \to 0$. Then use it to show that $\sin(1/x)$
is not uniformly continuous on $(0,1)$ and that $x^2$ is not uniformly
continuous on $\R^1$.\kn{The forward direction is a substitution. For the
converse, negate uniform continuity: some $\varepsilon$ has, for each
$\delta = 1/n$, a pair $p_n,q_n$ within $1/n$ whose images are
$\varepsilon$ apart --- the failing pair builds itself. For the examples take
$1/(2n\pi)$ against $1/(2n\pi + \pi/2)$, and $n$ against $n+1/n$.}
\provenance{Tao I, Prop.\ 9.9.8; ECNU \cjk{例}{Ex} 10--11.}

\item Prove the gluing lemma S4.16, and use it with Theorem 4.19 to show
that $\sqrt x$ is uniformly continuous on $[0,+\infty)$. Show likewise that
$\cos\sqrt x$ is uniformly continuous there.\kn{Split at $1$: on $[0,1]$ use
compactness, on $[1,\infty)$ use $|\sqrt x - \sqrt y| = |x-y|/(\sqrt x +
\sqrt y) \le \tfrac12|x-y|$. Note that $\sqrt x$ is uniformly continuous
although its derivative is unbounded near $0$ --- so ``bounded derivative''
is sufficient for uniform continuity but far from necessary.}
\provenance{ECNU \cjk{例}{Ex} 12 and \cjk{习题}{Ex} 4.2 \#12, \#20.}

\item Let $f$ be continuous on $[a,+\infty)$ and suppose
$\lim_{x\to+\infty} f(x)$ exists and is finite. Prove that $f$ is bounded on
$[a,+\infty)$ and uniformly continuous there. Must $f$ attain a maximum?
\kn{Both halves are the gluing pattern: the limit hypothesis makes $f$
nearly constant past some $M$, and $[a,M]$ is compact. So a hypothesis at
infinity substitutes for compactness of the domain. The last question is
no: $f(x) = 1 - 1/x$ on $[1,\infty)$ has supremum $1$, never attained.}
\provenance{ECNU \cjk{习题}{Ex} 4.2 \#6, \#16.}

\item Prove that if $f$ and $g$ are uniformly continuous and \emph{bounded}
on $E$, then $fg$ is uniformly continuous on $E$; and show by example that
boundedness cannot be dropped. Which of $f+g$, $\max(f,g)$, $g \circ f$ are
uniformly continuous without extra hypotheses?\kn{Split the product:
$|f(p)g(p) - f(q)g(q)| \le |f(p)|\,|g(p)-g(q)| + |g(q)|\,|f(p)-f(q)|$ --- the
identical decomposition as at 3.3(c). For the counterexample $f = g = x$ on
$\R^1$. Sum, max and composition all survive with no extra hypotheses.}
\provenance{Tao II, Ex.\ 2.3.4--2.3.6.}

\item Prove S4.17: a uniformly continuous $f : E \to \R^k$ has a limit at
every point of $\overline E$, and therefore extends uniquely to a
uniformly continuous function on $\overline E$. Show that $\sin(1/x)$ on
$(0,1)$ has no such extension, and identify which hypothesis it
violates.\kn{Uniform continuity sends Cauchy sequences to Cauchy sequences
(Rudin's Exercise 4.11), completeness converts those to limits, and S4.1
promotes ``every sequence converges'' to ``the limit exists.'' Uniqueness is
Rudin's Exercise 4.4. The example is merely continuous, not uniformly ---
by S4.15 with $p_n, q_n$ as in exercise 17.}
\provenance{Tao I, Cor.\ 9.9.14; generalises Rudin's Exercise 4.13.}

% ---------------- Group V ----------------
\item Define the oscillation $\omega_f(p)$ as in S4.4. Prove that $f$ is
continuous at $p$ if and only if $\omega_f(p) = 0$, that $\{x : \omega_f(x)
\ge \varepsilon\}$ is closed for each $\varepsilon > 0$, and hence that the
set of discontinuities of \emph{any} $f$ is a countable union of closed
sets. Deduce, using Exercise 3.22, that no function is continuous at every
rational and discontinuous at every irrational.\kn{The last part is the
converse question to Rudin's Exercise 4.18 (Thomae), and the answer is that
the two cases are not symmetric. If such an $f$ existed, $\Q$ would be a
countable intersection of dense open sets, and intersecting further with the
complements of the individual rationals would exhibit $\varnothing$ as such
an intersection --- contradicting Baire.}
\provenance{Yu Pin, Lemmas 131--132; Tao I, Ex.\ 9.3.4.}

\item Let $A \subset \R^1$ be countable. Construct a monotonically
increasing $f : \R^1 \to \R^1$ whose set of discontinuities is exactly
$A$.\kn{The exact converse of 4.30, and Rudin's Remark 4.31 sketches the
construction: enumerate $A = \{a_1,a_2,\dots\}$ and put $f(x) =
\sum_{a_n < x} 2^{-n}$. Convergence is comparison with a geometric series;
the jump at $a_n$ is exactly $2^{-n}$; continuity elsewhere follows because
the tail beyond $N$ contributes less than $2^{-N}$. Together with 4.30 this
characterises the possible discontinuity sets of monotone functions
completely --- a rare instance of a sharp answer.}
\provenance{Yu Pin, HW4 problem E (p.\,110); cf.\ Tao I, Ex.\ 9.8.5.}

\item $\star$~\emph{(Cayley--Hamilton by continuity.)} For $A$ an $n \times n$ complex
matrix let $P_A(X) = \det(XI - A)$. Prove: (a) if $A$ is diagonalisable then
$P_A(A) = 0$; (b) the map $A \mapsto P_A(A)$ is continuous on the space of
matrices with the norm $\|\cdot\|$; (c) every $A$ is a limit of
diagonalisable matrices; (d) therefore $P_A(A) = 0$ for every $A$.\kn{The
best advertisement in this chapter for Rudin's Exercise 4.4. Part (a) is
one line in an eigenbasis; (b) holds because the entries of $P_A(A)$ are
polynomials in the entries of $A$; (c) perturbs the eigenvalues of an upper
triangular form to make them distinct. Then a purely algebraic theorem falls
out of ``two continuous maps agreeing on a dense set are equal'' --- with
the density supplied by linear algebra and the rigidity by Chapter 4.}
\provenance{Yu Pin, winter problem set U (pp.\,331--332).}

\item Prove that a path-connected set is connected, and that the closure of
\[ S = \bigl\{\,(x,\sin(1/x)) : 0 < x \le 1\,\bigr\} \]
is connected but not path-connected.\kn{The first part is 4.22 applied to
the path. For the second, $\overline S$ adds the segment $\{0\}\times[-1,1]$;
connectedness holds because the closure of a connected set is connected,
while a path from the segment into $S$ would have to traverse infinitely
many oscillations in finite time --- make that precise by showing its first
coordinate cannot leave $0$. Rudin's favourite pathology at 4.27(d) turns
out to separate two notions of connectedness he never distinguishes.}
\provenance{Tao II, Ex.\ 2.4.7--2.4.9.}

\end{enumerate}
